% =============================================================================
% Chapter 6 — The Network Interface Controller ($A100–$A13F)
% =============================================================================
\chapter{The Network Interface Controller}

\epigraph{``The network is the computer.  The packet is the thought.''}%
         {\textit{--- Sun Microsystems proverb, paraphrased}}

\section*{Introduction}

The Network Interface Controller --- NIC for short --- provides message-oriented
TCP networking through four independent connection slots.  Each slot may operate
as a \emph{client} (connecting out to a remote host) or as a \emph{server}
(listening for incoming connections and accepting them).

Messages on the wire are length-prefixed: a single byte gives the payload
length, where a length byte of zero denotes 256 bytes.  This simple framing
removes the need for delimiter parsing on the 6502 side.

Data moves between the NIC and CPU RAM via DMA.  To send, the program writes
the source address and length into the DMA registers and issues a Send command.
To receive, the program issues a Recv command and the NIC copies the next
queued message into RAM at the DMA address.  Up to 16 messages may be queued
per slot before older messages are dropped.

The NIC occupies 64 bytes of address space from \$A100 to \$A13F.


% ═══════════════════════════════════════════════════════════════════════════════
\section{Command Register}
\label{sec:nic-cmd}
% ═══════════════════════════════════════════════════════════════════════════════

\mementry{A100}{NicCmd}{W}{---}

Writing a command byte to this register executes the command immediately on the
slot selected by \addr{A102}.  The available commands are:

\cmdentry{01}{Connect}{Connect to remote host}

Opens a TCP connection to the hostname in the name buffer (\addr{A120}) on the
port specified by \addr{A108}--\addr{A109}.  The connection is established
asynchronously; poll \addr{A118} (slot status) for the Connected bit, or check
the Error bit on failure.  A 10-second timeout applies.

\cmdentry{02}{Disconnect}{Close connection}

Closes the active connection on the selected slot and drains any unread
messages from the receive queue.

\cmdentry{03}{Send}{Send message via DMA}

Reads \addr{A112} bytes from CPU RAM starting at the DMA address
(\addr{A110}--\addr{A111}) and sends them as a single length-prefixed message
on the selected slot's TCP connection.  If the slot is not connected, the
slot's Error flag is set.

\cmdentry{04}{Recv}{Receive message via DMA}

Dequeues the next message from the selected slot's receive queue and writes it
into CPU RAM at the DMA address.  The length of the received message is stored
in \addr{A113}.  If no message is available, \addr{A113} reads as zero.

\cmdentry{05}{Listen}{Start listening for connections}

Begins listening for incoming TCP connections on the local port specified by
\addr{A10A}--\addr{A10B}.  The listener binds to the loopback interface.  When
a remote client connects, the slot's Data Ready bit is set in the slot status
register, indicating that an Accept command should be issued.

\cmdentry{06}{Accept}{Accept pending connection}

Accepts a pending inbound connection on the selected slot.  After acceptance,
the slot transitions to Connected and incoming messages are queued automatically.

\begin{notebox}
All commands operate on the slot selected by \addr{A102}.  Set the slot
\emph{before} writing any command, port, hostname, or DMA parameters.
\end{notebox}


% ═══════════════════════════════════════════════════════════════════════════════
\section{Status and Control Registers}
\label{sec:nic-status}
% ═══════════════════════════════════════════════════════════════════════════════

% --- $A101 NicStatus ----------------------------------------------------------

\mementry{A101}{NicStatus}{R}{---}

Global NIC status register.  Provides a quick summary of the controller state.

\bitfield{Err}{---}{---}{---}{---}{---}{Any}{Rdy}

\begin{description}
  \item[Bit 0 --- Ready]\hfill\\
    Always set to 1 once the NIC is initialised.
  \item[Bit 1 --- Any Data]\hfill\\
    Set if \emph{any} slot has data waiting in its receive queue (or a pending
    accept).  Useful for a single polling check across all slots.
  \item[Bit 7 --- Error]\hfill\\
    Set if any slot has encountered an error.
\end{description}

\begin{lstlisting}
10 REM wait for data on any slot
20 WAIT $A101, 2             : REM wait for bit 1 set
\end{lstlisting}

% --- $A102 NicSlot ------------------------------------------------------------

\mementry{A102}{NicSlot}{R/W}{0}

Active slot selector (0--3).  All commands, port settings, hostname writes, and
DMA operations act on the slot selected here.

\begin{lstlisting}
POKE $A102, 2 : REM select slot 2
\end{lstlisting}

% --- $A103 NicIrqCtrl ---------------------------------------------------------

\mementry{A103}{NicIrqCtrl}{R/W}{0}

IRQ enable bitmask.  Bit~$N$ (0--3) enables an IRQ when data arrives on
slot~$N$.  When a message is received on an enabled slot, the NIC asserts the
CPU's IRQ line.

\bitfield{---}{---}{---}{---}{Sl3}{Sl2}{Sl1}{Sl0}

\begin{lstlisting}
POKE $A103, 1   : REM enable IRQ on slot 0 only
POKE $A103, $0F : REM enable IRQ on all 4 slots
\end{lstlisting}

% --- $A104 NicIrqStatus -------------------------------------------------------

\mementry{A104}{NicIrqStatus}{R}{0}

IRQ pending flags.  Bit~$N$ is set when data has arrived on slot~$N$ (and IRQ
was enabled for that slot).  \textbf{Reading this register clears all pending
flags} and de-asserts the IRQ line.

\bitfield{---}{---}{---}{---}{Sl3}{Sl2}{Sl1}{Sl0}

\begin{warningbox}
Because reading clears the flags, read this register only once per interrupt
handler and save the value for inspection.
\end{warningbox}


% ═══════════════════════════════════════════════════════════════════════════════
\section{Port Registers}
\label{sec:nic-ports}
% ═══════════════════════════════════════════════════════════════════════════════

% --- $A108-$A109 NicRemotePort ------------------------------------------------

\memrange{A108}{A109}{NicRemotePort}{R/W}{---}

Remote port for the Connect command (16-bit, little-endian).  Write the low
byte to \addr{A108} and the high byte to \addr{A109}.

\begin{lstlisting}
DOKE $A108, 8080 : REM port 8080
\end{lstlisting}

% --- $A10A-$A10B NicLocalPort -------------------------------------------------

\memrange{A10A}{A10B}{NicLocalPort}{R/W}{---}

Local port for the Listen command (16-bit, little-endian).  Write the low byte
to \addr{A10A} and the high byte to \addr{A10B}.

\begin{lstlisting}
DOKE $A10A, 9000 : REM listen on port 9000
\end{lstlisting}


% ═══════════════════════════════════════════════════════════════════════════════
\section{DMA Registers}
\label{sec:nic-dma}
% ═══════════════════════════════════════════════════════════════════════════════

% --- $A110-$A111 NicDmaAddr ---------------------------------------------------

\memrange{A110}{A111}{NicDmaAddr}{R/W}{---}

CPU RAM address for DMA transfers (16-bit, little-endian).  For Send, this is
the source address.  For Recv, this is the destination address.

\begin{lstlisting}
DOKE $A110, $4000 : REM DMA to/from $4000
\end{lstlisting}

% --- $A112 NicDmaLen ----------------------------------------------------------

\mementry{A112}{NicDmaLen}{R/W}{0}

DMA transfer length for the Send command.  Valid values are 1--255.  A value of
0 is treated as 256 (the maximum message size).

% --- $A113 NicMsgLen ----------------------------------------------------------

\mementry{A113}{NicMsgLen}{R}{0}

Length of the last received message, set after a Recv command.  If no message
was available, this reads as 0.  A received message of exactly 256 bytes is
reported as 0 (matching the wire protocol convention).


% ═══════════════════════════════════════════════════════════════════════════════
\section{Per-Slot Status Registers}
\label{sec:nic-slot-status}
% ═══════════════════════════════════════════════════════════════════════════════

% --- $A118-$A11B NicSlotStatus ------------------------------------------------

\memrange{A118}{A11B}{NicSlotStatus0--3}{R}{---}

Four consecutive bytes, one per slot.  Each provides the current state of that
connection.  Slot~$N$'s status is at \$A118~+~$N$.

\bitfield{---}{---}{---}{RmCl}{Err}{SndR}{DatR}{Conn}

\begin{description}
  \item[Bit 0 --- Connected]\hfill\\
    Set when the slot has an active TCP connection.
  \item[Bit 1 --- Data Ready]\hfill\\
    Set when one or more messages are waiting in the receive queue, or when
    a pending accept is available (server mode).
  \item[Bit 2 --- Send Ready]\hfill\\
    Always set; indicates the slot can accept a Send command.
  \item[Bit 3 --- Error]\hfill\\
    Set on connection failure, send error, or receive queue overflow.
  \item[Bit 4 --- Remote Closed]\hfill\\
    Set when the remote peer has closed the connection.
\end{description}

\begin{lstlisting}
10 REM check if slot 0 is connected
20 IF PEEK($A118) AND 1 THEN PRINT "CONNECTED"
\end{lstlisting}

\begin{notebox}
Use \texttt{WAIT \$A118, 1} to block until slot~0 becomes connected.  The
\texttt{WAIT} instruction polls the address until the specified bits are set.
\end{notebox}


% ═══════════════════════════════════════════════════════════════════════════════
\section{Hostname Buffer}
\label{sec:nic-hostname}
% ═══════════════════════════════════════════════════════════════════════════════

% --- $A120-$A13F NicNameBuf ---------------------------------------------------

\memrange{A120}{A13F}{NicNameBuf}{R/W}{---}

A 32-byte buffer for the hostname used by the Connect command.  The name must
be null-terminated ASCII.  Write the hostname before issuing a Connect command.

\begin{lstlisting}
10 REM write "localhost" into the name buffer
20 A$ = "localhost" + CHR$(0)
30 FOR I=1 TO LEN(A$)
40   POKE $A120+I-1, ASC(MID$(A$,I,1))
50 NEXT I
\end{lstlisting}


% ═══════════════════════════════════════════════════════════════════════════════
\section{Complete Examples}
\label{sec:nic-examples}
% ═══════════════════════════════════════════════════════════════════════════════

\subsection{Client: Connect, Send, and Receive}

This example connects to a server on localhost port~8080, sends a
``HELLO'' message, and waits for a response.

\begin{lstlisting}
100 REM === TCP client example ===
110 POKE $A102, 0              : REM select slot 0
120 DOKE $A108, 8080           : REM remote port 8080
130 REM write hostname
140 A$ = "localhost" + CHR$(0)
150 FOR I=1 TO LEN(A$)
160   POKE $A120+I-1, ASC(MID$(A$,I,1))
170 NEXT I
180 POKE $A100, 1              : REM connect
190 WAIT $A118, 1              : REM wait for connected
200 PRINT "CONNECTED!"
210 REM
220 REM send "HELLO"
230 A$ = "HELLO"
240 FOR I=1 TO LEN(A$)
250   POKE $3FFF+I, ASC(MID$(A$,I,1))
260 NEXT I
270 DOKE $A110, $4000          : REM DMA source address
280 POKE $A112, LEN(A$)        : REM DMA length
290 POKE $A100, 3              : REM send
300 PRINT "SENT: "; A$
310 REM
320 REM receive response
330 DOKE $A110, $4100          : REM DMA dest address
340 WAIT $A118, 2              : REM wait for data ready
350 POKE $A100, 4              : REM recv
360 L = PEEK($A113)            : REM message length
370 R$ = ""
380 FOR I=0 TO L-1
390   R$ = R$ + CHR$(PEEK($4100+I))
400 NEXT I
410 PRINT "RECEIVED: "; R$
420 POKE $A100, 2              : REM disconnect
\end{lstlisting}

\subsection{Server: Listen, Accept, and Echo}

This example listens on port~9000, accepts one connection, and echoes back
every message it receives.

\begin{lstlisting}
100 REM === TCP echo server ===
110 POKE $A102, 0              : REM select slot 0
120 DOKE $A10A, 9000           : REM local port 9000
130 POKE $A100, 5              : REM listen
140 PRINT "LISTENING ON PORT 9000..."
150 WAIT $A118, 2              : REM wait for pending accept
160 POKE $A100, 6              : REM accept
170 PRINT "CLIENT CONNECTED!"
180 REM
190 REM echo loop
200 S = PEEK($A118)
210 IF S AND 16 THEN GOTO 310  : REM remote closed?
220 IF NOT(S AND 2) THEN 200   : REM no data? loop
230 DOKE $A110, $4000          : REM DMA address
240 POKE $A100, 4              : REM recv
250 L = PEEK($A113)
260 IF L=0 THEN 200            : REM empty message
270 POKE $A112, L              : REM set send length
280 POKE $A100, 3              : REM echo it back
290 PRINT "ECHOED "; L; " BYTES"
300 GOTO 200
310 PRINT "CLIENT DISCONNECTED"
320 POKE $A100, 2              : REM disconnect
\end{lstlisting}

\begin{tipbox}
The BASIC commands \texttt{NOPEN}, \texttt{NCLOSE}, \texttt{NLISTEN},
\texttt{NACCEPT}, and \texttt{NSEND} provide higher-level access to the NIC
without direct register manipulation.  See the NovaVM BASIC User Guide for
details.
\end{tipbox}
